% Appendix

% ----- A. Testnet Deployment Details -----
\subsection{Testnet Deployment Details}
\label{app:testnet}

\Cref{tab:testnet_txs} lists the transaction hashes for all on-chain operations
executed on the Polygon Amoy testnet (Chain ID: 80002).
All transactions are publicly verifiable on Polygonscan.\footnote{%
\url{https://amoy.polygonscan.com/address/0x2cC31B9EFBD5dDfB5456EE4c76A0fcBCA5EEE117}}

\begin{table}[ht]
\centering
\caption{Testnet transaction log on Polygon Amoy.}
\label{tab:testnet_txs}
\footnotesize
\setlength{\tabcolsep}{2pt}
\begin{tabular}{@{}llr@{}}
\toprule
\textbf{Operation} & \textbf{TX Hash (truncated)} & \textbf{Gas} \\
\midrule
Deploy
  & \texttt{0xb230...6bf5}
  & 1{,}824{,}117 \\
\texttt{mint()} (Tanaka)
  & \texttt{0x5461...8c30}
  & 189{,}792 \\
\texttt{mintBatch(10)}
  & \texttt{0xfcee...9182}
  & 1{,}025{,}449 \\
\texttt{revoke()}
  & \texttt{0x9fad...5ecd}
  & 50{,}027 \\
\bottomrule
\end{tabular}
\end{table}

Contract address: \hexid{0x2cC31B9EFBD5dDfB5456EE4c76A0fcBCA5EEE117}\\
Deployer address: \hexid{0x91eca41e060de63871dFb0Fe8Fdc78897d6119b8}\\
Deployment block: 35{,}058{,}633


% ----- B. Unit Test Results -----
\subsection{Unit Test Results}
\label{app:tests}

\Cref{tab:test_results} lists the v1 suite, developed with Hardhat against the
deployed contract (solc 0.8.20, optimizer 200 runs). The further 35 tests summarised
in \cref{subsubsec:v2} cover co-signing by a studio attester and its idempotence
across multiple studios, refusal to co-sign a revoked token, attester-role
management, ERC-5192 signalling, order-independence and value-sensitivity of the
canonical serialisation, the disclosure commitment of \cref{subsec:privacy}, and one
end-to-end run from metadata construction through mint and co-signature to
verification.

\begin{table}[ht]
\centering
\caption{Unit test suite for the deployed v1 \texttt{AnimatorSBT} contract
(35 tests, all passing). A further 35 cover the v2 co-signing reference contract,
canonical serialisation, disclosure commitments and a local end-to-end flow,
for 70 in total.}
\label{tab:test_results}
\footnotesize
\setlength{\tabcolsep}{2pt}
\begin{tabular}{@{}ll@{}}
\toprule
\textbf{Category} & \textbf{Test Case} \\
\midrule
\multirow{4}{*}{Deployment}
  & Set correct name and symbol \\
  & Grant DEFAULT\_ADMIN\_ROLE to deployer \\
  & Grant ISSUER\_ROLE to deployer \\
  & Start with 0 total minted \\
\midrule
\multirow{5}{*}{\texttt{mint()}}
  & Mint with ISSUER\_ROLE \\
  & Emit SBTMinted event \\
  & Record issuedAt timestamp \\
  & Increment token IDs \\
  & Revert when caller lacks ISSUER\_ROLE \\
\midrule
\multirow{4}{*}{\texttt{mintBatch()}}
  & Batch mint to multiple recipients \\
  & Batch mint 10 tokens (gas measurement) \\
  & Revert on length mismatch \\
  & Revert when caller lacks ISSUER\_ROLE \\
\midrule
\multirow{2}{*}{\texttt{tokenURI()}}
  & Return correct URI after mint \\
  & Revert for non-existent token \\
\midrule
\multirow{4}{*}{Soulbinding}
  & Revert on \texttt{transferFrom} \\
  & Revert on \texttt{safeTransferFrom} \\
  & Revert on \texttt{approve} \\
  & Revert on \texttt{setApprovalForAll} \\
\midrule
\multirow{5}{*}{\texttt{revoke()}}
  & Revoke a token \\
  & Emit SBTRevoked event \\
  & Revert on double revoke \\
  & Revert for non-existent token \\
  & Revert when caller lacks ISSUER\_ROLE \\
\midrule
\multirow{6}{*}{View Functions}
  & \texttt{isRevoked()} returns false for active token \\
  & \texttt{issuedAt()} returns block timestamp \\
  & \texttt{totalMinted()} returns correct count \\
  & \texttt{tokensOfOwner()} returns all tokens \\
  & \texttt{tokensOfOwner()} returns single token \\
  & \texttt{tokensOfOwner()} returns empty for non-holder \\
\midrule
\multirow{4}{*}{AccessControl}
  & Admin grants ISSUER\_ROLE \\
  & Admin revokes ISSUER\_ROLE \\
  & Non-admin cannot grant roles \\
  & \texttt{supportsInterface} (ERC721, AccessControl) \\
\midrule
Gas & Measure deploy gas \\
\bottomrule
\end{tabular}
\end{table}

% ----- C. Metadata JSON Example -----
\subsection{Metadata JSON Example}
\label{app:metadata}

\Cref{lst:metadata_json} shows the actual metadata JSON uploaded to IPFS
during the E2E demonstration (Tanaka scenario).
The IPFS CID is \hexid{Qmbv6Xm1GmNpxkJ8rq3QCihTPmQYBo6EQaxE72qTcXtVys},
retrievable at \url{https://gateway.pinata.cloud/ipfs/Qmbv6Xm1GmNpxkJ8rq3QCihTPmQYBo6EQaxE72qTcXtVys}.

\begin{lstlisting}[
  basicstyle=\scriptsize\ttfamily,
  frame=single,
  caption={SBT metadata JSON from E2E demo (Tanaka scenario).},
  label={lst:metadata_json},
  breaklines=true,
  columns=flexible
]
{
  "name": "AnimatorSBT - Tanaka Yuki",
  "description": "Contribution certification
    for Tanaka Yuki on Episode 7",
  "attributes": {
    "animator_name": "Tanaka Yuki",
    "animator_wallet": "0x7eBa20e703...82f3a",
    "project": "Episode 7",
    "role": "key_anim",
    "task_count": 25,
    "start_date": "2025-01-15",
    "end_date": "2025-03-30",
    "studio_name": "Studio Example",
    "issued_at": "2026-03-11T16:04:05.123Z"
  },
  "evidence_files": [
    "ipfs://QmExampleHash1_storyboard",
    "ipfs://QmExampleHash2_keyframes"
  ],
  "evidence_hash": "0x7c27a833f1678308
    c58d217c178a1c27ac10a32b4a8b61ec7d681dbe24673e79"
}
\end{lstlisting}

The \texttt{evidence\_hash} here is the Keccak-256 hash of the \texttt{attributes}
object as serialised, and any party holding the metadata can recompute it; what that
does and does not establish is set out in \cref{subsec:verification}.

This file is also the record of the defect \cref{subsec:privacy} corrects:
\texttt{animator\_name} sits in cleartext on public storage, and the digest is taken
over an object containing it with no salt. We print the deployed artifact rather than
a regenerated one because it is what the recorded transactions point at.

% ----- D. Smart Contract Source Code -----
\subsection{Smart Contract Source Code}
\label{app:source}

The \texttt{AnimatorSBT} contract is 160 lines of Solidity and is not reproduced
here; it is in the repository named under \emph{Data and code availability}, together
with the v2 co-signing reference implementation, the issuance service, the
verification portal, the test suite and the measurement and cost scripts. Two
properties of the deployed version matter for reading the rest of this appendix. It
is written against OpenZeppelin Contracts v4, so soulbinding is enforced by
overriding \texttt{\_beforeTokenTransfer} and token identifiers come from the
\texttt{Counters} utility; v5 removes both, and a production port would enforce the
same rule in the unified \texttt{\_update} hook, as noted in \cref{subsec:stack}.
And \texttt{revoke()} sets a flag, leaving the token and its URI in place, which is
the behaviour \cref{subsec:privacy} examines.

% ----- E. Co-signing Additions (v2) -----
\subsection{Co-signing Additions (AnimatorSBTV2)}
\label{app:v2source}

\Cref{lst:v2_cosign} reproduces the core additions of the v2 reference contract
(\cref{subsubsec:v2}): the studio-attester role, the idempotent multi-studio
\texttt{coSign}, the trust-level view, and ERC-5192 soulbound signaling. The
remainder of the contract is identical to v1. As stated in
\cref{subsec:cosigning}, this contract is validated by the test suite of
\cref{app:tests} but is \emph{not} deployed on-chain, and in the single-operator
MVP the issuer also holds \texttt{STUDIO\_ATTESTER\_ROLE}, so a co-signature adds no
authenticity until an independent studio operates the role.

\begin{lstlisting}[
  basicstyle=\scriptsize\ttfamily,
  frame=single,
  caption={AnimatorSBTV2 --- co-signing and ERC-5192 additions (excerpt).},
  label={lst:v2_cosign},
  breaklines=true,
  columns=flexible,
  morekeywords={function,external,public,view,override,returns,return,if,revert,
    mapping,event,error,emit,address,uint256,uint8,bool,bytes4,memory,calldata,
    constant,onlyRole,indexed}
]
bytes32 public constant STUDIO_ATTESTER_ROLE =
  keccak256("STUDIO_ATTESTER_ROLE");
bytes4 private constant _INTERFACE_ID_ERC5192 =
  0xb45a3c0e;

mapping(uint256 => address[]) private _attesters;
mapping(uint256 => mapping(address => bool))
  private _hasCoSigned;

event SBTCoSigned(uint256 indexed tokenId,
  address indexed attester, uint256 timestamp);
event Locked(uint256 tokenId); // ERC-5192

error CannotCoSignRevokedToken(uint256 tokenId);
error DuplicateAttestation(
  uint256 tokenId, address attester);

// Studio (msg.sender) co-signs an existing token.
function coSign(uint256 tokenId)
  external onlyRole(STUDIO_ATTESTER_ROLE) {
  _requireMinted(tokenId);
  if (_revoked[tokenId])
    revert CannotCoSignRevokedToken(tokenId);
  if (_hasCoSigned[tokenId][msg.sender])
    revert DuplicateAttestation(tokenId, msg.sender);
  _hasCoSigned[tokenId][msg.sender] = true;
  _attesters[tokenId].push(msg.sender);
  emit SBTCoSigned(
    tokenId, msg.sender, block.timestamp);
}

// 0 = self-attested, 1 = >=1 studio co-signature.
function attestationLevel(uint256 tokenId)
  external view returns (uint8) {
  return _attesters[tokenId].length > 0 ? 1 : 0;
}

// ERC-5192: every minted token is permanently locked.
function locked(uint256 tokenId)
  external view returns (bool) {
  _requireMinted(tokenId);
  return true;
}
\end{lstlisting}
